\originalchapter{来源与改编}\label{ch:source}
\originalsection{逐讲对应}
本书主源是 Casey Rodriguez 授课、Andrew Lin 记录的 MIT 18.102 Spring 2021. 合订本125个 PDF 物理页, 23份独立讲义合计147个物理页; 差别来自分讲排版及前后信息, 并不是另外147页独立课程内容. 官方目录重复链接第1讲, 已按实际下载文件去重. 下表中页码为125页合订本的物理页序; 讲次交界页可能同时包含前讲末尾和后讲日期, 因此边界按下一讲日期所在页划分, 只用作定位.

\begin{longtable}{p{10mm}p{20mm}p{125mm}}
\toprule
讲次 & 合订页序 & 内容及本书对应\\
\midrule
\endhead
1 & 1--6 & 向量空间、范数、连续有界函数和 Banach 空间: 第\ref{ch:norm}章.\\
2 & 7--13 & 绝对级数、有界算子、算子空间和对偶: 第\ref{ch:norm}、\ref{ch:operators}、\ref{ch:dual}章.\\
3 & 14--17 & 子空间、商空间、Baire 与统一有界: 第\ref{ch:operators}、\ref{ch:baire}章.\\
4 & 18--21 & 开映射、闭图及 Zorn 引理: 第\ref{ch:baire}、\ref{ch:dual}章.\\
5 & 22--25 & Hamel 基与 Hahn--Banach: 第\ref{ch:dual}章, 补全复数保范数证明.\\
6 & 26--30 & 双对偶及外测度: 双对偶在第\ref{ch:dual}章完整展开; 外测度在第\ref{ch:measure}章给定义接口.\\
7 & 31--35 & $\sigma$ 代数、可测集和 Borel 集: 第\ref{ch:measure}章定义接口; 构造证明交由实变科.\\
8 & 36--40 & Lebesgue 测度的构造与性质: 第\ref{ch:measure}章明示调用, 未逐项重编构造.\\
9 & 41--46 & 可测函数及其运算: 第\ref{ch:measure}章必要定义与闭包性质.\\
10 & 47--52 & 简单函数逼近及非负积分: 第\ref{ch:measure}章.\\
11 & 53--58 & 单调收敛、Fatou: 第\ref{ch:measure}章完整证明.\\
12 & 59--64 & 可积函数、控制收敛和 Riemann 积分关系: 第\ref{ch:measure}章完整证明控制收敛及连续函数的两种积分一致性.\\
13 & 65--70 & $L^p$、Hölder、Minkowski、完备性、内积: 第\ref{ch:measure}、\ref{ch:hilbert}章完整证明.\\
14 & 71--76 & Hilbert 空间、正交系、Gram--Schmidt、Bessel: 第\ref{ch:hilbert}章.\\
15 & 77--82 & 正交基、Parseval、可分分类、Fourier 与 Dirichlet 核: 第\ref{ch:hilbert}、\ref{ch:fourier}章.\\
16 & 83--88 & Fejér 核与 Fourier 完整性: 第\ref{ch:fourier}章, 图形按定义重绘.\\
17 & 89--93 & 闭凸集极小化、正交分解、Riesz 表示: 第\ref{ch:adjoint}章.\\
18 & 94--99 & 伴随、紧集与一致小尾项: 第\ref{ch:adjoint}、\ref{ch:compact}章.\\
19 & 100--104 & 紧集判别、有限秩算子: 第\ref{ch:compact}章.\\
20 & 105--109 & 紧算子、谱、Neumann 级数、自伴范数: 第\ref{ch:compact}、\ref{ch:spectrum}章.\\
21 & 110--113 & 自伴谱端点、紧自伴特征空间: 第\ref{ch:spectrum}章, 补全端点向量范数估计.\\
22 & 114--118 & Fredholm 条件、最大原则、紧自伴谱定理: 第\ref{ch:spectrum}章.\\
23 & 119--125 & Green 核、正弦谱、平方根、Dirichlet 问题: 第\ref{ch:dirichlet}章, 校正核符号并补全正则性.\\
\bottomrule
\end{longtable}

\originalsection{补充、勘误与边界}
编者补充包括赋范空间完备化及稠密延拓, 数列空间对偶的计算证明, 平行四边形恒等式的逆向证明, 一般弱与弱星拓扑、可分对偶球弱星紧性, 及配套例题、习题的完整解答. 第\ref{ch:weak}章的 Hilbert 弱收敛与第\ref{ch:compact}章的紧算子关系另参 Richard Melrose 2020讲义第3章第15节, 印刷页82--84. 这些补充均围绕本册的依赖关系, 不把辅助讲义的所有章节纳入主源范围.

原第4讲的一处说明把标准坐标向量称作 $\ell^1$ 的 Hamel 基; 它们的有限张成实际是 $c_{00}$, 稠密却不是全空间. 原第5讲对复数 Hahn--Banach 的实虚部相容性略写, 本册改用实化与单位相位估计. 原第13讲某个衰减条件的不等号写成下界, 其证明使用的是上界; 本册的积分例子以直接积分给出准确条件. 原第16讲 Fejér 核质量陈述的积分区间误写为整条实轴, 而周期核在实轴上的积分不是1; 本册使用一个周期的积分. 该讲的半角公式也按正确的 $1-\cos t=2\sin^2(t/2)$ 使用.

原第21讲从二次型趋于零到 $\norm{(A-\lambda I)x_n}\to0$ 的推理需要额外估计, 已在自伴谱端点证明中补全. 原第22讲的谱基先在非闭值域上讨论, 本册明确把完整 Hilbert 展开放在值域闭包 $(\ker A)^\perp$ 中. 原第23讲的分段 Green 核符号与 $-u''=f$ 不一致, 已以正核 $\min(x,y)(1-\max(x,y))$ 修正; 非零特征值为 $(k\pi)^{-2}$, 平方根的上确界估计依赖 $\norm{f}_2$, 且乘法算子只需有界, 不应被默认称为紧算子. 本书通过已证明的正弦基、统一收敛及 Green 表示补齐这些论证.

本册完整覆盖主源的 Banach 与 Hilbert 空间主线、Fourier 完备性、紧算子、自伴谱及连续非负势的区间 Dirichlet 问题, 并收入官方全部10份作业、期中测验和期末作业的题面及完整解答. 原第6--12讲全部测度构造和非可测集细节仍未逐页重编, 但考核实际要求的外测度平移与缩放、开集区间分解、有限区间并逼近、Littlewood 三原则、$L^\infty$ 完备性与 $L^p$ 可分性均在对应题解内证明.

未覆盖视频文字稿, Carleson 定理的证明, 一般不可分 Banach--Alaoglu 定理, 一般 Banach 空间自反性的等价刻画, 非自伴 Fredholm 全理论, 一般非紧自伴谱测度、无界算子的域理论、一般 Sobolev 与分布、迹类和 Schatten 理论. 原作业指定的周期 Fourier 加权 $H^s$ 空间、$s>1/2$ 连续嵌入、$s>0$ 紧嵌入及磨光逼近已收入, 不能再把这几个具体结果列作未覆盖, 也不能由此宣称覆盖一般 Sobolev 理论. 更广专题有些出现在176页 Melrose 辅助讲义中, 不由参考书目录推断本册完成范围.

\originalsection{作业与考核对应}
下表中的页数为各原PDF的物理页数, 包含最后的 MIT OCW 署名页. “单元”指显式字母子问, 或没有字母分问的原编号单题. 期末第2题的两个积分按一个原编号单元登记, 解答中逐个计算. 完整原题条件、子问及中文定位另见 assessment-inventory.json; 原件没有官方答案, 本书的新增解答统一标为 AI 独立解答.
\begin{longtable}{p{32mm}r r r p{63mm}}
\toprule
材料 & 原页数 & 题数 & 单元 & 本书定位\\
\midrule\endhead
作业1 & 3 & 5 & 9 & 第\ref{sec:ps01}节, Hölder、完备、数列对偶.\\
作业2 & 3 & 5 & 10 & 第\ref{sec:ps02}节, 可逆性、商空间、闭值域和反例.\\
作业3 & 3 & 4 & 9 & 第\ref{sec:ps03}节, 对偶转置、外测度、开集.\\
作业4 & 2 & 3 & 4 & 第\ref{sec:ps04}节, 可测性与集合逼近.\\
作业5 & 3 & 4 & 9 & 第\ref{sec:ps05}节, 扩展实值运算与 Egorov.\\
作业6 & 3 & 4 & 7 & 第\ref{sec:ps06}节, 连续逼近、Riemann--Lebesgue.\\
作业7 & 2 & 4 & 9 & 第\ref{sec:ps07}节, 可分性、$L^\infty$、极化.\\
作业8 & 3 & 4 & 9 & 第\ref{sec:ps08}节, 正交、$H^s$ 与连续嵌入.\\
作业9 & 4 & 5 & 10 & 第\ref{sec:ps09}节, 磨光、凸集投影.\\
作业10 & 4 & 6 & 15 & 第\ref{sec:ps10}节, 紧性、移位谱、Volterra.\\*
期中测验 & 3 & 5 & 8 & 第\ref{ch:midterm}章.\\*
期末作业 & 3 & 5 & 6 & 第\ref{ch:final}章.\\
\bottomrule
\end{longtable}
12份合计36个原PDF物理页、54道编号题、105个解答单元, 即87个字母子问与18道未分问单题.

原12章的30份例题与习题解答全部保留, 仅增交叉引用标签. 直接对应的复用包括: 题\ref{ps03:1}复用例题\ref{legacy:ch03:example1}的延拓泛函; 题\ref{ps07:1}(a)复用习题\ref{legacy:ch05:exercise1}的有限测度包含; 题\ref{midterm:3}(b)复用习题\ref{legacy:ch04:exercise1}的完备范数比较; 题\ref{ps09:3}复用习题\ref{legacy:ch07:exercise1}的自伴幂等投影; 题\ref{final:4}(b)复用习题\ref{legacy:ch10:exercise1}的对角紧性. 题\ref{ps08:1}(b)把习题\ref{legacy:ch06:exercise1}的有限正交近似推广到闭包, 题\ref{final:5}细化例题\ref{legacy:ch11:example1}的乘法谱; 都保留所需扩展论证, 不以一个章节号代替新增目标的证明. 其余题目对既有定理的精确引用均在解答内列明.

源题勘误还包括作业2第4题同构的箭头, 作业5递减测度连续性的名称, 作业8第1题漏印范数, 作业9前置子列定理的字母与阈值混写. 作业9第4题缺少非空假设使原命题按字面为假, 本书明确给出空集反例再补条件. 作业8第2题与期末第3题的闭包横线按原PDF保留. 所有这些修订均在题面附近说明.

原件目录保留实际下载的 PDF、TeX 与原源码 ZIP, 每个文件的字节、SHA256、页数和许可依据均可在源清单定位. 本书中文源码 ZIP 是独立可编译的干净源码, 不包含字体、编译缓存、页面图像或原 ZIP 内的 macOS 元数据. 原源归档与中文成册的范围在两个目录的说明中分别列明.
